\documentclass{article}
\usepackage{graphicx} % Required for inserting images
\usepackage[T1]{fontenc}
\usepackage[margin=1in]{geometry}
\usepackage{booktabs}
\usepackage{tabularx}
\usepackage[hidelinks]{hyperref}

\title{Using an Agent to Decide Which Algo to Run}
\author{Strategy research \textperiodcentered{} L/S Gamma Desk, QUANTT}
\date{2026-09-27}

\begin{document}

\maketitle

\textbf{Question:} look into how an agent can be used to decide which algo to run. this should be very similar to the agents are not algorithms paper

\textbf{Method:} Searched all 22 papers in \texttt{docs/\allowbreak{}papers/\allowbreak{}} (page-tagged index) + web search. Quotes verified against the cited PDF pages. Page numbers are PDF pages; printed page numbers are given where they differ.

\textbf{Related research:} \href{2026-09-27-deep-hedging-as-gamma-scalping-pnl-driver.pdf}{Deep hedging as the gamma scalping P\&L driver} (covers OPHR's hedger-routing agent)

\section{Short answer}

\emph{Agents Are Not Algorithms} (Cheng, Granger, Shi \& Strela 2026) gives a clear design rule: use the LLM agent only for the sparse \textbf{judgment} decision, feed it \textbf{pre-computed numbers}, and hand everything \textbf{time-sensitive to compiled code}. Done that way, their agent became competitive with a fixed-rule algorithm. When the agent also ran the fast execution loop, it did much worse.

For the desk this maps directly:

\begin{itemize}

\item \textbf{The agent's job:} the once-a-day choice of which \texttt{LG\_\allowbreak{}}/\allowbreak{}\texttt{SG\_\allowbreak{}} algo to run, or to stay flat.

\item \textbf{Its inputs:} the pipeline's features (RV forecast, VRP, VIX term structure, VVIX, trend, ATM IV), calculated beforehand.

\item \textbf{Compiled code, never the agent:} strike selection, order execution (\texttt{broker.\allowbreak{}py}) and intraday hedging (\texttt{Hedging/\allowbreak{}}).

\end{itemize}

The paper does \textbf{not} show that an agent beats a well-designed rule. The rule stayed about as good. So the agent must be tested against a simple compiled router. Separately, an LLM can't be backtested fairly on history it may have seen during training, so forward paper trading is the real test.

\section{What the papers say}

\subsection{Agents reason when the decision is made, and that takes time}

\begin{itemize}

\item \textbf{What makes an agent different from an algorithm:}

\begin{quote}``Agentic AI systems built on large language models reason about each decision in real time rather than executing a pre-specified policy.'' (Cheng, Granger, Shi \& Strela 2026, \emph{Agents Are Not Algorithms}, p. 1)\end{quote}

\item \textbf{An RL policy trained in advance counts as an algorithm, not an agent:}

\begin{quote}``Algorithms, whether rule-based or learned through reinforcement learning before deployment, represent the fully compiled policy endpoint.'' (Cheng et al., p. 2, printed p. 1)\end{quote}

\item \textbf{The paper doesn't claim agents win:}

\begin{quote}``To be clear, our question is not whether agents or algorithms dominate, but how and where decision-time reasoning is valuable.'' (Cheng et al., p. 2, printed p. 1)\end{quote}

\end{itemize}

\subsection{The task splits into judgment and speed}

\begin{itemize}

\item \textbf{Their task has the same two parts as the desk's:}

\begin{quote}``Tender acceptance rewards judgment: the trader must evaluate quoted edge, displayed liquidity, remaining time, inventory, and penalties. Execution rewards speed'' (Cheng et al., p. 3, printed p. 2)\end{quote}

For them, ``tender acceptance'' means deciding whether to take a client's block trade. For the desk, the equivalent judgment is which algo to run today.

\item \textbf{Reasoning helps the decision but hurts execution:}

\begin{quote}``Reasoning improves the agent's intended decisions and the quality of trades it does complete, but delay creates a wedge between intention and realized trading performance.'' (Cheng et al., p. 4, printed p. 3)\end{quote}

\item \textbf{Slow decisions are where reasoning pays:}

\begin{quote}``In our slowest treatment, the same agent earns the highest average profits.'' (Cheng et al., p. 4, printed p. 3)\end{quote}

\item \textbf{More reasoning isn't automatically better:}

\begin{quote}``The robust feature is the non-monotonicity: increasing decision-time reasoning does not mechanically improve trading performance.'' (Cheng et al., p. 25, printed p. 24)\end{quote}

The best setting depends on the model:

\begin{quote}``For the Anthropic models we test, performance is highest with extended thinking turned off and declines as the reasoning budget rises.'' (Cheng et al., p. 6, printed p. 5)\end{quote}

\end{itemize}

\subsection{The architecture that worked}

\begin{itemize}

\item \textbf{Pre-computed decision support, with no recommendation attached:}

\begin{quote}``This decision-support block reports book-based tender edge, projected inventory, limit-breach indicators, and flattening costs, but does not recommend an action.'' (Cheng et al., p. 22, printed p. 21)\end{quote}

\begin{quote}``In the low-reasoning cell, mean NLV rises from \$18.6 thousand for the baseline agent to \$30.4 thousand with pre-computed calculations.'' (Cheng et al., p. 23, printed p. 22)\end{quote}

NLV is net liquidation value, i.e. the account's ending value.

\item \textbf{The agent only makes the judgment call; compiled code executes:}

\begin{quote}``The reasoning-off agent with compiled unwind earns a slightly lower mean NLV of \$44.5 thousand, is profitable in every session, and has the highest Sharpe ratio in this architecture, 2.85.'' (Cheng et al., p. 23, printed p. 22)\end{quote}

\begin{quote}``These outcomes are competitive with a fully compiled algorithm that uses a 5-cent threshold to accept tenders, which earns mean NLV of \$44.5 thousand with a Sharpe ratio of 2.49.'' (Cheng et al., p. 23, printed p. 22)\end{quote}

\item \textbf{The design rule in one sentence:}

\begin{quote}``calculations and low-latency execution can be compiled into the surrounding system, while the agent is most valuable when reserved for the judgment margin.'' (Cheng et al., p. 26, printed p. 25)\end{quote}

\item \textbf{For their order-routing task, the fixed rule won:}

\begin{quote}``when the routing, tender-evaluation, and unwind logic can be specified in advance, the algorithm avoids much of the real-time deliberation cost faced by the agent.'' (Cheng et al., p. 74, printed p. 73)\end{quote}

\end{itemize}

\subsection{An RL agent picking the strategy (OPHR)}

\begin{itemize}

\item \textbf{Two agents: one picks the position, one picks the hedger:}

\begin{quote}``Option Position Agent (OP-Agent) controlling the long/\allowbreak{}short of volatility, and ii) Hedger Routing Agent (HR-Agent) selecting the optimal Hedgers with different risk preferences to perform dynamic Delta hedging based on positions and market conditions.'' (Chen, Cai, Qin \& An, \emph{OPHR}, p. 2)\end{quote}

\item \textbf{The router picks, and the chosen hedger executes:}

\begin{quote}``The HR-Agent makes a decision every N steps, and the selected Hedger executes hedging at each step.'' (\emph{OPHR}, p. 5)\end{quote}

\item \textbf{Both share one goal:}

\begin{quote}``Both agents share the common objective of maximizing portfolio net value while managing risk exposure.'' (\emph{OPHR}, p. 4)\end{quote}

\item \textbf{Training starts from an imperfect rule, then improves on it:}

\begin{quote}``we first distill a sub-optimal Oracle policy to OP-Agent, then alternatively train HR-Agent and OP-Agent.'' (\emph{OPHR}, p. 2)\end{quote}

\end{itemize}

\subsection{A fixed-rule router (PRISM)}

\begin{itemize}

\item \textbf{PRISM routes with fixed ``gates'' instead of an agent:}

\begin{quote}``The VRC investment committee requires that all four of the following gates pass before the VRC derivatives desk may initiate any long gamma position.'' (Verma 2026, \emph{PRISM}, p. 34)\end{quote}

The four gates are a VRP threshold, an arbitrage condition, a minimum cost-adjusted Sharpe, and regime certainty.

\item \textbf{One gate blocks trades while the regime is unclear:}

\begin{quote}``This gate prevents VRC from entering a long gamma position during regime transition ambiguity'' (\emph{PRISM}, p. 35)\end{quote}

PRISM is a non-peer-reviewed whitepaper and its results are simulated.

\end{itemize}

\section{How the papers connect}

\begin{center}\small
\begin{tabularx}{\linewidth}{>{\raggedright\arraybackslash}X >{\raggedright\arraybackslash}X >{\raggedright\arraybackslash}X >{\raggedright\arraybackslash}X}
\toprule
\textbf{Paper} & \textbf{Builds on /\allowbreak{} agrees with} & \textbf{Differs from} & \textbf{Key contribution} \\
\midrule
Cheng et al. 2026, \emph{Agents Are Not Algorithms} & Agrawal, Gans \& Goldfarb (prediction vs judgment); Lopez-Lira 2025 & Treats RL policies as compiled algorithms & Place the agent only on the judgment decision; pre-compute its inputs; compile execution \\
OPHR & Murray et al. 2022; Buehler et al. & An RL router, trained before deployment, i.e. ``compiled'' in Cheng's terms & Separate agents for vol position and hedger choice \\
PRISM & Whalley--Wilmott; HAR-RV & No learning; fixed gates & A rule-based router: VRP, regime and cost gates \\
\emph{The Alpha Illusion} (web) & --- & End-to-end LLM trading papers & Reported LLM alpha is not deployment evidence \\
\emph{Agentic Trading} survey (web) & --- & --- & 77 studies; weak evaluation standards \\
\bottomrule
\end{tabularx}
\end{center}

\begin{itemize}

\item \textbf{Three designs, same shape.} Cheng's agent-plus-compiled-execution, OPHR's routing agent over fixed hedgers, and PRISM's gate rules all separate \emph{choosing} from \emph{executing}. They differ in whether the choice is made by reasoning at the time (LLM), a trained policy (RL), or fixed rules.

\item \textbf{The rule is always a strong baseline.} Cheng's fixed rule matched the best agent setup and beat the agent on the routing task. PRISM runs entirely on rules.

\item \textbf{Pre-computed inputs matter.} Cheng's decision-support gain matches OPHR giving its router Greeks and position data as inputs.

\end{itemize}

\section{Relevance to the LS Gamma desk}

These are suggestions, not decisions the desk has made:

\begin{itemize}

\item \textbf{The desk's current structure already matches Cheng's winning architecture.}

\begin{center}\small
\begin{tabularx}{\linewidth}{>{\raggedright\arraybackslash}X >{\raggedright\arraybackslash}X}
\toprule
\textbf{Cheng et al.} & \textbf{LS Gamma desk} \\
\midrule
Tender judgment (agent) & Daily choice of which \texttt{LG\_\allowbreak{}}/\allowbreak{}\texttt{SG\_\allowbreak{}} algo to run, or none (the future signals/\allowbreak{}agent layer) \\
Decision support (pre-computed, no recommendation) & \texttt{pipeline/\allowbreak{}} output: \texttt{features} (RV/\allowbreak{}HAR, trend, VIX term structure, VIX percentile, VVIX, rf), \texttt{atm\_\allowbreak{}iv}, plus the RV forecast and VRP \\
Compiled unwind execution & \texttt{broker.\allowbreak{}submit()}, and the algo files' fixed legs \\
Continuous trading loop kept away from the agent & \texttt{Hedging/\allowbreak{}} hedgers running on a schedule, with no LLM in that loop \\
\bottomrule
\end{tabularx}
\end{center}

\item \textbf{A daily decision is the ``slow market'' case,} where Cheng et al. found reasoning pays off most (p. 4). That's an inference: their fastest decisions were seconds, the desk's is daily.

\item \textbf{Constrain the agent's output.} Suggestion: it returns only an algo name from a fixed menu (the registry of \texttt{LG\_\allowbreak{}}/\allowbreak{}\texttt{SG\_\allowbreak{}} algos plus ``stay flat'') and a short reason. Code checks that choice before anything runs. The agent never calls \texttt{broker} directly, and never picks strikes unless that's added deliberately later.

\item \textbf{Test against three routers:}

\begin{enumerate}

\item \textbf{A compiled rule router}, PRISM-style gates. For example, ``VRP above threshold \ensuremath{\rightarrow} short gamma structure X; term structure inverted \ensuremath{\rightarrow} iron condor'', with thresholds you set. This is the baseline that must be beaten.

\item \textbf{An LLM agent} with the pipeline features as input, choosing from the menu. Test its reasoning setting (off /\allowbreak{} low /\allowbreak{} higher), because the best level depends on the model (p. 6, p. 25).

\item \textbf{Optionally, an RL router} like OPHR's, which is compiled in Cheng's terms. It needs historical training data, e.g. the club's R2 dataset.

\end{enumerate}

\item \textbf{Log every decision:} the inputs, the menu choice and the agent's stated reason. Cheng et al. studied recorded reasoning (p. 4) and it makes the router auditable for the desk. \texttt{research/\allowbreak{}Agentic\_\allowbreak{}Routing/\allowbreak{}} is the natural home for these experiments.

\end{itemize}

\section{Gaps and open questions}

\begin{itemize}

\item \textbf{Different task.} Cheng et al. study block-trade acceptance and inventory unwinding in a simulator (Rotman Interactive Trader), over seconds and minutes. They don't test picking between options strategies, or daily decisions.

\item \textbf{Agent \ensuremath{\approx} rule, not agent > rule.} In their best setup the agent was only \emph{competitive} with a 5-cent threshold rule. Whether an agent adds value over a simple VRP/\allowbreak{}regime rule for the desk is untested.

\item \textbf{Backtesting an LLM is contaminated.} An LLM may have seen historical market outcomes in its training data (\emph{Alpha Illusion}, web). Running it over 2014--2025 history isn't a fair test, so forward paper trading is the credible evidence. An RL or rule router doesn't have this problem.

\item \textbf{Model- and prompt-dependent.} Cheng et al. find the best reasoning level differs across models and providers. The prompt, the menu design and the input format all need testing.

\item \textbf{Where strike selection lives.} The papers don't say whether choosing strikes and expiries is ``judgment'' or should be compiled. Starting compiled is the safer default.

\end{itemize}

\section{Sources}

\textbf{Repo papers used, with pages cited:}

\begin{itemize}

\item \texttt{Agents-\allowbreak{}Are-\allowbreak{}not-\allowbreak{}Algorithms.\allowbreak{}pdf}: Cheng, Granger, Shi \& Strela 2026, \emph{Agents Are Not Algorithms: The Tradeoffs of Decision-Time Reasoning in AI Trading} (pp. 1, 2, 3, 4, 5, 6, 22, 23, 25, 26, 74)

\item \texttt{OPHR-\allowbreak{}Mastering-\allowbreak{}Volatility-\allowbreak{}Trading-\allowbreak{}with-\allowbreak{}Multi-\allowbreak{}Agent-\allowbreak{}Deep-\allowbreak{}Reinforcement-\allowbreak{}Learning.\allowbreak{}pdf}: Chen, Cai, Qin \& An (pp. 2, 4, 5)

\item \texttt{PRISM-\allowbreak{}A-\allowbreak{}Probabilistic-\allowbreak{}Regime-\allowbreak{}Integrated-\allowbreak{}Scalping-\allowbreak{}Model-\allowbreak{}for-\allowbreak{}Derivatives-\allowbreak{}Arbitrage.\allowbreak{}pdf}: Verma 2026 (pp. 34, 35)

\end{itemize}

\textbf{Checked, nothing on choosing between strategies:} the other 19 papers. The deep hedging and RL papers choose hedge sizes, not strategies. The forecasting and VRP papers provide the inputs, not the routing.

\textbf{Web sources (outside the repo):}

\begin{itemize}

\item \href{https://papers.ssrn.com/abstract=6713620}{Cheng et al. 2026 on SSRN}, \href{https://inghawcheng.github.io/}{author page}, \href{https://mlquants.substack.com/p/agents-are-not-algorithms-the-tradeoffs}{ML Quants summary}

\item \href{https://arxiv.org/abs/2605.16895}{Ye et al. 2026, \emph{The Alpha Illusion: Reported Alpha from LLM Trading Agents Should Not Be Treated as Deployment Evidence}}

\item \href{https://arxiv.org/abs/2605.19337}{Xia et al. 2026, \emph{Agentic Trading: When LLM Agents Meet Financial Markets} (survey of 77 studies)}

\end{itemize}

\textbf{Suggested papers to add to \texttt{docs/\allowbreak{}papers/\allowbreak{}}:}

\begin{enumerate}

\item \textbf{The Alpha Illusion (2026):} how to evaluate an LLM router without fooling yourself; minimum reporting standards.

\item \textbf{Agentic Trading survey (2026):} a map of the LLM-agent designs that have been tried, and which studies meet basic evidence standards.

\item \textbf{Murray et al. 2022:} the hedger family OPHR's router chooses from; useful if the desk builds a routing agent over \texttt{Hedging/\allowbreak{}}.

\item \textbf{Agrawal, Gans \& Goldfarb, \emph{Prediction, Judgment, and Complexity}:} the prediction-vs-judgment framing Cheng et al. build on. Found only through their references, not on the web.

\end{enumerate}

\end{document}
